\begin{quote}\small Historical Round28 checkpoint. References below to planned, unexecuted or next work describe the state at that checkpoint; the integrated Round29 chapter supersedes that plan.\end{quote}
\section{Completed investigation count and remaining obligations}
\label{r28:sec:closeout}
All ten requested investigations are admitted, completing five adaptive goal pairs. This is 100\% of the requested investigation count, including limited outcomes. It is not a percentage of the Yang--Mills existence-and-mass-gap proof: equations, loops, numerical error percentages and network size provide no calibrated denominator for that theorem.

The final assessment identifies the strongest structural advance as a controlled, nonvacuous physical sector of the declared fixed-spacing SU(2) homogeneous omitted interaction on its selected-strip reference. AJ1 supplies the normal local representation and reducing physical generator; AJ2 supplies a nonzero original Wilson fluctuation; AK1/AK2 add conditional variance, physical form-energy, finite-window mass and imaginary-time control. These are improvements within the specified whole-star orthant state. The inherited $\alpha/16$ threshold is not a new continuum gap, and the unnormalized Wilson spectral mass is not a particle count or an evaluated mass.

The final retrospective assessment is one contemporary model-agent review applying the documented historical-method questions. It is not five new independent reviewers, participation by historical figures, external human peer review or another scientific investigation. The independent productions share inherited premises and model authorship. The Newton AK2 comparison is by the reverse producer; the separate skeptical production and later clarifications retain their recorded attribution.

The advisor's 22 September 2026 \href{https://www.claymath.org/millennium/yang-mills-the-maths-gap/}{official problem-status check} records the problem as unsolved. Selected sections 3--5 and 6.5 of the \href{https://www.claymath.org/wp-content/uploads/2022/06/yangmills.pdf}{Jaffe--Witten description} were checked for scope, not as an exhaustive current research survey. The required nontrivial four-dimensional theory for any compact simple gauge group, suitable axioms, required short-distance behavior and strictly positive finite physical mass gap remain open here. Confinement and an isolated one-particle state are discussed extensions, not additional mandatory clauses of the stated Millennium problem. Numerical lattice stability constants, a particular correction iteration and exact Wilson moment equality are possible technical milestones, likewise not compulsory clauses in themselves.

For this workbench's proposed route, full generator, observable, state, coupling, spacing and clock matching remains unresolved. So does identification of specified bulk and boundary limits. A boundary-local orthant witness cannot silently become a bulk observable, and the finite, canonical, summable and homogeneous models cannot be interchanged. The finite first-moment equality has not been proved to survive the limit; no actual strict moment loss is established.

The final roadmap ranks the following \emph{unselected} questions. Its ordering concerns missing implications and bounded test value, not success probabilities or equal units of proof work.
\begin{center}\small
\begin{tabularx}{\textwidth}{@{}p{10mm}p{43mm}X@{}}
\toprule Goal & Proposed direction & First unresolved premise \\
\midrule
AL & Full model and physical-scale bridge & A complete coefficient, observable and clock dictionary with one admissible scaling path; compare the existing bare-coupling obstruction before claiming a new bridge.\\[4pt]
AM & Numerical homogeneous stability interval & Evaluate every constant required for the complete original model, retaining an insufficient verdict if a necessary constant or infinite estimate remains open.\\[4pt]
AN & Specified bulk-boundary comparison & Compare deep-bulk orthant translates with one declared centered-box route using a bound that vanishes with boundary distance.\\[4pt]
AO & Same-Wilson moment identification & Prove sufficient uniform integrability or form convergence before transferring the finite identity as equality; assume no new energy-tail or second-moment theorem.\\[4pt]
AP & Canonical inference at a prior-selected design & Separately test one design fixed from coarse prior information and decisions across independently varying candidate parameters.\\
\bottomrule
\end{tabularx}
\end{center}
Every second loop remains unselected until its first review. A certified numerical evaluator for the declared 561-state AH enrichment is a lower-priority computational alternative; no such heat vector was produced by this cycle's assessment. The detailed models, dependency IDs and proposed first tests are recorded in \repo{research/round28/advisor/roadmap.json}. Investigation ten is the stopping boundary. This closeout executes no further research target.

The assessment, roadmap and official scope check are separately source-bound. They do not alter the 33-record exploratory reading ledger or claim new primary research by the retrospective reviewer. Scientific priority remains unverified; the two-star marker identifies contributions in this workbench.
